% Generated by task tex-catalogue from dossiers/code-layer1.md, sections 0 and A.
% Snippets are the dossier's, copied from the source at b435631 (bound proofs elided where the dossier says so).
\section{Layer 1, declaration by declaration}\label{sec:gen-cat-layer1}

\subsection{The objects Layer 1 is built on}

\textbf{CompPoly's hypercube table} (CompPoly is the Verified-zkEVM library of computable polynomials; revision \texttt{3468b38c}, \texttt{CompPoly/\allowbreak{}Multilinear/\allowbreak{}Basic.\allowbreak{}lean}). A multilinear polynomial in \texttt{n} variables is stored by its \texttt{2\textasciicircum{}n} values on the Boolean cube, as a vector; index \texttt{i} is the cube point whose coordinate \texttt{k} is bit \texttt{k} of \texttt{i} (low bit first):

\begin{leancode}
-- CompPoly/Multilinear/Basic.lean:45-47 at 3468b38c
  The indexing is **little-endian** (i.e. the least significant bit is the first bit). -/
@[reducible]
def CMlPolynomialEval (R : Type*) (n : ℕ) := Vector R (2 ^ n) -- coefficient of Lagrange basis
\end{leancode}

\texttt{lagrangeBasis w} is the vector of the equality kernel \texttt{eq(\allowbreak{}w,\allowbreak{} x) = ∏\_\allowbreak{}k (\allowbreak{}x\_\allowbreak{}k w\_\allowbreak{}k + (\allowbreak{}1-{}x\_\allowbreak{}k)(\allowbreak{}1-{}w\_\allowbreak{}k))} over the cube points \texttt{x}; \texttt{evalMle t x} is the multilinear extension of the table \texttt{t} evaluated at a point \texttt{x ∈ R\textasciicircum{}n} (by folding one variable at a time, low variable first); \texttt{eval₂Mle t φ x} first maps the entries through a ring homomorphism \texttt{φ : R →+* S} (for leanVM, the embedding \texttt{algebraMap K E} of \texttt{K = GF(\allowbreak{}2\textasciicircum{}64)} into \texttt{E = GF(\allowbreak{}2\textasciicircum{}192)}), then evaluates at \texttt{x ∈ S\textasciicircum{}n}:

\begin{leancode}
-- CompPoly/Multilinear/Basic.lean:410-411, 499-500, 520-521 at 3468b38c
def lagrangeBasis (w : Vector R n) : Vector R (2 ^ n) :=
  Vector.ofFn (fun i => ∏ j : Fin n, if (BitVec.ofFin i).getLsb j then w[j] else 1 - w[j])
def evalMle (p : CMlPolynomialEval R n) (x : Vector R n) : R :=
  evalMleValues p x
def eval₂Mle (p : CMlPolynomialEval R n) (f : R →+* S) (x : Vector S n) : S :=
  evalMle (map f p) x
\end{leancode}

\textbf{The spine's objects Layer 1 plugs into} (\texttt{LeanerVM/\allowbreak{}Protocol/\allowbreak{}Field.\allowbreak{}lean}, \texttt{Spine/\allowbreak{}Instance.\allowbreak{}lean}, \texttt{Spine/\allowbreak{}Seams.\allowbreak{}lean} at \texttt{b435631}). \texttt{Column n} wraps a \texttt{K}-table of \texttt{n} variables; its oracle answers a query \texttt{r ∈ E\textasciicircum{}n} with \texttt{eval₂Mle q (\allowbreak{}algebraMap K E) r} (\texttt{Field.\allowbreak{}lean:102-{}110}, \texttt{evalOracle\_\allowbreak{}answer} is \texttt{rfl}). A \texttt{Layout μ ι κ} is a \emph{reading law}: how to read column \texttt{c} (height \texttt{2\textasciicircum{}κ c}) off a stack of \texttt{2\textasciicircum{}μ} cells, how to lift a point of the column to a point of the stack, and the law that the two agree:

\begin{leancode}
-- LeanerVM/Protocol/Spine/Instance.lean:73-81
structure Layout (μ : ℕ) (ι : Type) (κ : ι → ℕ) where
  /-- Read column `c` off the stack. -/
  read : Column μ → (c : ι) → Column (κ c)
  /-- Lift a point of column `c` to the point of the stack that reads the same value. -/
  extend : (c : ι) → Vector E (κ c) → Vector E μ
  /-- The selector law: reading then extending is extending then reading. -/
  read_eval : ∀ (q : Column μ) (c : ι) (z : Vector E (κ c)),
    CMlPolynomialEval.eval₂Mle (read q c).values (algebraMap K E) z =
      CMlPolynomialEval.eval₂Mle q.values (algebraMap K E) (extend c z)
\end{leancode}

An \texttt{M3Instance} carries one such \texttt{layout} for all its columns (\texttt{Instance.\allowbreak{}lean:140}), and the relation \texttt{M3Holds} reads every value through it. A \texttt{Weight μ} is a weight on the stack given by its cube values \texttt{onCube} and an evaluator \texttt{mle} with the proof \texttt{mle\_\allowbreak{}eq} that it computes the extension of \texttt{onCube}; \texttt{Weight.\allowbreak{}pair W q = Σ\_\allowbreak{}i W.\allowbreak{}onCube[i] · q[i]}; a \texttt{ColumnClaim} is (column, point, value) and holds when the column's extension at the point is the value; a \texttt{WeightedClaim} is (weight, value) and holds when the pairing is the value (\texttt{Seams.\allowbreak{}lean:57-{}67,\allowbreak{} 103-{}124}).

\textbf{Mathlib vocabulary.} \texttt{Antitone f} means \texttt{a ≤ b → f b ≤ f a} (non-increasing). \texttt{Vector R n} is an array of length \texttt{n}; \texttt{z ++ s} concatenates, so \texttt{z} fills the low coordinates. \texttt{Fin m} is \texttt{\{0,\allowbreak{} …,\allowbreak{} m-{}1\}}.

\subsection{Classes, and how this catalogue is laid out}

Class legend. \textbf{L} load-bearing: a definition the instance \texttt{leanIsaInstance} (Layer 3) or the executable verifier (Layer 12) will unfold, so it becomes trusted surface when they are built (no master theorem, seam or built instance unfolds any Layer 1 declaration today; see E). \textbf{I} interface: consumed by name by a later phase, the adaptor, or listed in the blueprint's interface list. \textbf{H} helper: a proof step. \textbf{T} test fixture (none: fixtures live under \texttt{tests/}). Counts are public declarations (the two \texttt{private} lemmas are not counted).

For each module: the load-bearing definitions as the dossier copied them, with their meaning in words; then a table of the load-bearing and interface declarations; then a table of the helpers. The dossier gives a snippet for the load-bearing definitions only; the interface declarations are listed by name, line and meaning.

\subsection{\file{LeanerVM/Protocol/ToCompPoly/Multilinear.lean}}

(generic: no protocol vocabulary; 43 public.)

Load-bearing definitions, copied:

\begin{leancode}
-- LeanerVM/Protocol/ToCompPoly/Multilinear.lean:144-149
def cubeIndex {k m : ℕ} (i : Fin (2 ^ k)) (j : Fin (2 ^ m)) : Fin (2 ^ (k + m)) :=
  ⟨i.val + 2 ^ k * j.val, by … ⟩            -- (proof of the bound elided here, 4 lines)
-- :170-171
def cubeSplit (k m : ℕ) : Fin (2 ^ k) × Fin (2 ^ m) ≃ Fin (2 ^ (k + m)) :=
  (Equiv.prodComm _ _).trans (finProdFinEquiv.trans (finCongr (by rw [pow_add, mul_comm])))
-- :235-236
def boolVec {m : ℕ} (j : Fin (2 ^ m)) : Vector R m :=
  Vector.ofFn fun b ↦ if j.val.testBit b then 1 else 0
-- :307-309
def slice {k m : ℕ} (t : CMlPolynomialEval R (k + m)) (j : Fin (2 ^ m)) :
    CMlPolynomialEval R k :=
  Vector.ofFn fun i ↦ t[cubeIndex i j]
\end{leancode}

\src{LeanerVM/Protocol/ToCompPoly/Multilinear.lean, at b435631; further ranges in the comments}

In words: \texttt{cubeIndex i j} is the index with low bits \texttt{i} and high bits \texttt{j}; \texttt{cubeSplit} is the same map as an equivalence (\texttt{cubeSplit\_\allowbreak{}apply} proves it equals \texttt{cubeIndex}); \texttt{boolVec j} is the cube point of index \texttt{j} (coordinate \texttt{b} is bit \texttt{b}); \texttt{slice t j} is the subtable of \texttt{t} whose high \texttt{m} bits are \texttt{j} (an aligned block).

\par\medskip\noindent\textbf{Load-bearing and interface declarations.}

{\footnotesize
\begin{longtable}{@{}>{\raggedright\arraybackslash}p{0.280\linewidth}>{\raggedright\arraybackslash}p{0.080\linewidth}>{\raggedright\arraybackslash}p{0.260\linewidth}>{\raggedright\arraybackslash}p{0.300\linewidth}@{}}
\toprule
Declaration & Line & Class & Meaning \\
\midrule
\endfirsthead
\toprule
Declaration & Line & Class & Meaning \\
\midrule
\endhead
\bottomrule
\endfoot
\texttt{sumCube t} & 67 & interface & \texttt{Σ\_\allowbreak{}i t[i]} \\
\texttt{hadamard s t} & 70 & interface (listed) & pointwise product \\
\texttt{evalMle\_\allowbreak{}eq\_\allowbreak{}sumCube\_\allowbreak{}hadamard} & 91 & interface (listed) & same, as \texttt{sumCube (\allowbreak{}hadamard …)} \\
\texttt{cubeIndex} & 144 & load-bearing & above \\
\texttt{cubeSplit}, \texttt{cubeSplit\_\allowbreak{}apply} & 170, 173 & load-bearing, helper & above \\
\texttt{lowVec}, \texttt{highVec} & 185, 189 & interface & low \texttt{k} / high \texttt{m} coordinates of a point (inside \texttt{lowPoint}, \texttt{highPoint}) \\
\texttt{boolVec} & 235 & load-bearing & above \\
\texttt{onesIndex m} & 239 & interface & the index \texttt{2\textasciicircum{}m − 1} (used by \texttt{padHigh}) \\
\texttt{boolVec\_\allowbreak{}zero} & 243 & interface & \texttt{boolVec 0 = (\allowbreak{}0,\allowbreak{}…,\allowbreak{}0)} (used by \texttt{PublicInput}) \\
\texttt{slice} & 307 & load-bearing & above \\
\texttt{evalMle\_\allowbreak{}append\_\allowbreak{}boolVec} & 342 & interface (listed; used by \texttt{PublicInput}) & a Boolean high part selects a slice \\
\texttt{placeSlice t j} & 351 & interface (listed) & \texttt{t} on the slice \texttt{j}, zero elsewhere (\texttt{padHigh}, \texttt{windowTable}) \\
\texttt{evalMle\_\allowbreak{}placeSlice} & 371 & interface (listed) & extension of a placed table = \texttt{(\allowbreak{}lagrangeBasis s)[j]·t̃(\allowbreak{}z)} \\
\end{longtable}}

\par\medskip\noindent\textbf{Helpers.}

{\footnotesize
\begin{longtable}{@{}>{\raggedright\arraybackslash}p{0.300\linewidth}>{\raggedright\arraybackslash}p{0.080\linewidth}>{\raggedright\arraybackslash}p{0.200\linewidth}>{\raggedright\arraybackslash}p{0.340\linewidth}@{}}
\toprule
Declaration & Line & Class & Meaning \\
\midrule
\endfirsthead
\toprule
Declaration & Line & Class & Meaning \\
\midrule
\endhead
\bottomrule
\endfoot
\texttt{hadamard\_\allowbreak{}getElem} & 73 & helper & entry of the product \\
\texttt{lagrangeBasis\_\allowbreak{}getElem\_\allowbreak{}nat} & 78 & helper & \texttt{lagrangeBasis\_\allowbreak{}getElem} with a \texttt{ℕ} index \\
\texttt{evalMle\_\allowbreak{}eq\_\allowbreak{}sum} & 85 & helper & \texttt{evalMle t x = Σ\_\allowbreak{}i t[i]·(\allowbreak{}lagrangeBasis x)[i]} \\
\texttt{evalMle\_\allowbreak{}cast}, \texttt{eval₂Mle\_\allowbreak{}cast} & 97, 103 & helper & evaluation invariant under a cast of \texttt{n} (used by \texttt{PublicInput}) \\
\texttt{sumCube\_\allowbreak{}lagrangeBasis} & 125 & helper & partition of unity \\
\texttt{evalMle\_\allowbreak{}replicate} & 131 & helper & a constant table extends to the constant \\
\texttt{cubeIndex\_\allowbreak{}val}, \texttt{cubeIndex\_\allowbreak{}div}, \texttt{cubeIndex\_\allowbreak{}mod}, \texttt{testBit\_\allowbreak{}cubeIndex} & 151-164 & helper & arithmetic of \texttt{cubeIndex} \\
\texttt{sum\_\allowbreak{}cube\_\allowbreak{}split} & 179 & helper & a cube sum is a double sum, high outside \\
\texttt{lowVec\_\allowbreak{}append}, \texttt{highVec\_\allowbreak{}append}, \texttt{lowVec\_\allowbreak{}append\_\allowbreak{}highVec} & 193-208 & helper & split and join are inverse \\
\texttt{lagrangeBasis\_\allowbreak{}cubeIndex} & 220 & helper & the basis factors across the index split \\
\texttt{boolVec\_\allowbreak{}onesIndex} & 250 & helper & \texttt{boolVec (\allowbreak{}2\textasciicircum{}m−1) = (\allowbreak{}1,\allowbreak{}…,\allowbreak{}1)} \\
\texttt{fin\_\allowbreak{}eq\_\allowbreak{}iff\_\allowbreak{}testBit} & 257 & helper & indices equal iff low bits equal \\
\texttt{lagrangeBasis\_\allowbreak{}boolVec} & 269 & helper & basis at a cube point is an indicator \\
\texttt{lagrangeBasis\_\allowbreak{}onesIndex}, \texttt{lagrangeBasis\_\allowbreak{}zero\_\allowbreak{}index} & 285, 292 & helper & basis at the two corners \\
\texttt{evalMle\_\allowbreak{}boolVec} & 298 & helper & the extension at a cube point reads the entry \\
\texttt{slice\_\allowbreak{}getElem}, \texttt{cubeIndex\_\allowbreak{}lt}, \texttt{slice\_\allowbreak{}getElem\_\allowbreak{}nat} & 312-323 & helper & entries of a slice \\
\texttt{evalMle\_\allowbreak{}split} & 331 & helper & \texttt{evalMle t (\allowbreak{}z ++ s) = Σ\_\allowbreak{}j (\allowbreak{}lagrangeBasis s)[j]·evalMle (\allowbreak{}slice t j) z} \\
\texttt{slice\_\allowbreak{}placeSlice} & 357 & helper & slicing a placed table \\
\texttt{sumCube\_\allowbreak{}placeSlice} & 381 & helper & placing keeps the cube sum \\
\end{longtable}}

\subsection{\file{LeanerVM/Protocol/ToCompPoly/BitProductTable.lean}}

(generic; 13 public.)


\begin{leancode}
-- LeanerVM/Protocol/ToCompPoly/BitProductTable.lean:127-128
def powersTable (a : R) (n : ℕ) : CMlPolynomialEval R n :=
  Vector.ofFn fun i ↦ a ^ i.val
\end{leancode}

\src{LeanerVM/Protocol/ToCompPoly/BitProductTable.lean, at b435631; further ranges in the comments}


\par\medskip\noindent\textbf{Load-bearing and interface declarations.}

{\footnotesize
\begin{longtable}{@{}>{\raggedright\arraybackslash}p{0.280\linewidth}>{\raggedright\arraybackslash}p{0.080\linewidth}>{\raggedright\arraybackslash}p{0.260\linewidth}>{\raggedright\arraybackslash}p{0.300\linewidth}@{}}
\toprule
Declaration & Line & Class & Meaning \\
\midrule
\endfirsthead
\toprule
Declaration & Line & Class & Meaning \\
\midrule
\endhead
\bottomrule
\endfoot
\texttt{bitProductTable f} & 44 & interface (listed) & entry \texttt{i} is \texttt{∏\_\allowbreak{}k f k (\allowbreak{}bit k of i)} \\
\texttt{evalMle\_\allowbreak{}bitProductTable} & 58 & interface (listed) & extension \texttt{= ∏\_\allowbreak{}k (\allowbreak{}(\allowbreak{}1 − x\_\allowbreak{}k) f k false + x\_\allowbreak{}k f k true)} \\
\texttt{powersTable a n} & 127 & load-bearing (inside \texttt{idxColumn}) & entry \texttt{i} is \texttt{a\textasciicircum{}i} \\
\texttt{evalMle\_\allowbreak{}powersTable} & 144 & interface (listed) & extension \texttt{= ∏\_\allowbreak{}k (\allowbreak{}(\allowbreak{}1 − x\_\allowbreak{}k) + x\_\allowbreak{}k a\textasciicircum{}\{2\textasciicircum{}k\})} \\
\texttt{evalMle\_\allowbreak{}lagrangeBasis} & 156 & interface (field \texttt{mle\_\allowbreak{}eq} of \texttt{eqWeight}) & the extension of \texttt{lagrangeBasis w} at \texttt{x} is \texttt{∏\_\allowbreak{}k (\allowbreak{}(\allowbreak{}1−x\_\allowbreak{}k)(\allowbreak{}1−w\_\allowbreak{}k) + x\_\allowbreak{}k w\_\allowbreak{}k)} \\
\end{longtable}}

\par\medskip\noindent\textbf{Helpers.}

{\footnotesize
\begin{longtable}{@{}>{\raggedright\arraybackslash}p{0.300\linewidth}>{\raggedright\arraybackslash}p{0.080\linewidth}>{\raggedright\arraybackslash}p{0.200\linewidth}>{\raggedright\arraybackslash}p{0.340\linewidth}@{}}
\toprule
Declaration & Line & Class & Meaning \\
\midrule
\endfirsthead
\toprule
Declaration & Line & Class & Meaning \\
\midrule
\endhead
\bottomrule
\endfoot
\texttt{bitProductTable\_\allowbreak{}getElem}, \texttt{map\_\allowbreak{}bitProductTable} & 47, 51 & helper & entries; commutes with a ring map \\
\texttt{eval₂Mle\_\allowbreak{}bitProductTable} & 86 & helper & same at a point of another ring \\
\texttt{lagrangeBasis\_\allowbreak{}eq\_\allowbreak{}bitProductTable} & 93 & helper & the Lagrange basis is such a table \\
\texttt{pow\_\allowbreak{}eq\_\allowbreak{}prod\_\allowbreak{}testBit} & 100 & helper & \texttt{a\textasciicircum{}i = ∏\_\allowbreak{}k (\allowbreak{}a\textasciicircum{}\{2\textasciicircum{}k\})\textasciicircum{}\{bit k of i\}} \\
\texttt{powersTable\_\allowbreak{}eq\_\allowbreak{}bitProductTable}, \texttt{map\_\allowbreak{}powersTable} & 130, 137 & helper &  \\
\texttt{eval₂Mle\_\allowbreak{}powersTable} & 149 & helper (proof of \texttt{idxColumn\_\allowbreak{}eval}) & same at a point of another ring \\
\end{longtable}}

\subsection{\file{LeanerVM/Protocol/ToCompPoly/Stacking.lean}}

(generic; 38 public.)


\begin{leancode}
-- LeanerVM/Protocol/ToCompPoly/Stacking.lean:62-68
structure Blocks where
  /-- The number of blocks. -/
  n : ℕ
  /-- The number of variables of each block; its height is `2 ^ size b`. -/
  size : Fin n → ℕ
  /-- Largest first. -/
  descending : Antitone size
-- :80-87
def offsetNat (k : ℕ) : ℕ :=
  ∑ c ∈ Finset.range k, if h : c < B.n then 2 ^ B.size ⟨c, h⟩ else 0
def offset (b : Fin B.n) : ℕ := B.offsetNat b.val
def total : ℕ := B.offsetNat B.n
-- :159-163
def stackAt (t : B.Tables R) (μ : ℕ) (pad : R) : CMlPolynomialEval R μ :=
  Vector.ofFn fun x ↦
    if B.total ≤ x.val then pad
    else ∑ b : Fin B.n,
      if B.InWindow b x.val then ((t b)[x.val - B.offset b]?).getD 0 else 0
-- :213-214 (bound proof elided, 6 lines)
def selector {μ : ℕ} (hμ : B.total ≤ 2 ^ μ) (b : Fin B.n) : Fin (2 ^ (μ - B.size b)) :=
  ⟨B.offset b / 2 ^ B.size b, by … ⟩
-- :227-230
def extendPoint {μ : ℕ} (hμ : B.total ≤ 2 ^ μ) (b : Fin B.n) (z : Vector R (B.size b)) :
    Vector R μ :=
  Vector.cast (Nat.add_sub_cancel' (B.size_le hμ b))
    (z ++ (boolVec (B.selector hμ b) : Vector R (μ - B.size b)))
-- :264-267
def unstack {μ : ℕ} (hμ : B.total ≤ 2 ^ μ) (q : CMlPolynomialEval R μ)
    (b : Fin B.n) : CMlPolynomialEval R (B.size b) :=
  slice (Vector.cast (congrArg (2 ^ ·) (Nat.add_sub_cancel' (B.size_le hμ b)).symm) q)
    (B.selector hμ b)
\end{leancode}

\src{LeanerVM/Protocol/ToCompPoly/Stacking.lean, at b435631; further ranges in the comments}

In words: a \texttt{Blocks} is a list of sizes, non-increasing; block \texttt{b} starts at \texttt{offset b}, the sum of the heights before it; \texttt{total} is the sum of all heights; \texttt{stackAt} lays given tables end to end and fills the rest with \texttt{pad}; \texttt{selector b = offset b /\allowbreak{} 2\textasciicircum{}size b}; \texttt{extendPoint b z = (\allowbreak{}z,\allowbreak{} bits of selector b)}; \texttt{unstack q b} reads block \texttt{b}'s window off any table \texttt{q}.

\par\medskip\noindent\textbf{Load-bearing and interface declarations.}

{\footnotesize
\begin{longtable}{@{}>{\raggedright\arraybackslash}p{0.280\linewidth}>{\raggedright\arraybackslash}p{0.080\linewidth}>{\raggedright\arraybackslash}p{0.260\linewidth}>{\raggedright\arraybackslash}p{0.300\linewidth}@{}}
\toprule
Declaration & Line & Class & Meaning \\
\midrule
\endfirsthead
\toprule
Declaration & Line & Class & Meaning \\
\midrule
\endhead
\bottomrule
\endfoot
\texttt{Blocks} & 62 & load-bearing & above \\
\texttt{Blocks.\allowbreak{}Tables R} & 75 & interface (listed) & one table per block \\
\texttt{offsetNat}, \texttt{offset}, \texttt{total} & 80, 84, 87 & load-bearing & above \\
\texttt{pow\_\allowbreak{}size\_\allowbreak{}dvd\_\allowbreak{}offset} & 109 & interface (sketch) & alignment: \texttt{2\textasciicircum{}size b ∣ offset b} \\
\texttt{stackAt} & 159 & interface (listed; the honest prover's stack) & above \\
\texttt{selector}, \texttt{selector\_\allowbreak{}val} & 213, 222 & load-bearing, helper & above \\
\texttt{extendPoint} & 227 & load-bearing & above \\
\texttt{lowPoint}, \texttt{highPoint} & 233, 238 & interface & the low \texttt{size b} / high \texttt{μ − size b} coordinates of a stack point \\
\texttt{unstack} & 264 & load-bearing & above \\
\texttt{unstack\_\allowbreak{}getElem} & 271 & interface (the cell meaning of the reader) & \texttt{(\allowbreak{}unstack q b)[i] = q[i + offset b]} \\
\texttt{unstack\_\allowbreak{}stackAt} & 280 & interface & reading the stack returns the table \\
\texttt{unstack\_\allowbreak{}eval₂} & 330 & interface (listed) & \texttt{eval₂Mle q φ (\allowbreak{}extendPoint b z) = eval₂Mle (\allowbreak{}unstack q b) φ z}, every \texttt{q} \\
\texttt{stack\_\allowbreak{}eval}, \texttt{stack\_\allowbreak{}eval₂} & 339, 345 & interface (listed) & the same on the stack, any pad \\
\end{longtable}}

\par\medskip\noindent\textbf{Helpers.}

{\footnotesize
\begin{longtable}{@{}>{\raggedright\arraybackslash}p{0.300\linewidth}>{\raggedright\arraybackslash}p{0.080\linewidth}>{\raggedright\arraybackslash}p{0.200\linewidth}>{\raggedright\arraybackslash}p{0.340\linewidth}@{}}
\toprule
Declaration & Line & Class & Meaning \\
\midrule
\endfirsthead
\toprule
Declaration & Line & Class & Meaning \\
\midrule
\endhead
\bottomrule
\endfoot
\texttt{offsetNat\_\allowbreak{}succ}, \texttt{offsetNat\_\allowbreak{}mono} & 89, 93 & helper &  \\
\texttt{offset\_\allowbreak{}add\_\allowbreak{}pow\_\allowbreak{}le\_\allowbreak{}offset}, \texttt{offset\_\allowbreak{}add\_\allowbreak{}pow\_\allowbreak{}le\_\allowbreak{}total} & 97, 103 & helper & windows disjoint, in order, inside the total \\
\texttt{InWindow b x} & 118 & helper & \texttt{offset b ≤ x < offset b + 2\textasciicircum{}size b} \\
\texttt{inWindow\_\allowbreak{}unique}, \texttt{exists\_\allowbreak{}inWindow}, \texttt{inWindow\_\allowbreak{}iff\_\allowbreak{}div} & 121, 129, 147 & helper & one window per index; no gaps below \texttt{total}; window = quotient \\
\texttt{stackAt\_\allowbreak{}getElem\_\allowbreak{}of\_\allowbreak{}inWindow}, \texttt{stackAt\_\allowbreak{}getElem\_\allowbreak{}of\_\allowbreak{}total\_\allowbreak{}le} & 165, 176 & helper & the cells of the stack \\
\texttt{map\_\allowbreak{}stackAt}, \texttt{stackAt\_\allowbreak{}eq\_\allowbreak{}of\_\allowbreak{}total\_\allowbreak{}eq} & 182, 197 & helper & commutes with ring maps; a full stack ignores the pad \\
\texttt{size\_\allowbreak{}le} & 206 & helper & every block fits in \texttt{μ} variables \\
\texttt{cast\_\allowbreak{}lowPoint\_\allowbreak{}append\_\allowbreak{}highPoint}, \texttt{lowPoint\_\allowbreak{}extendPoint}, \texttt{highPoint\_\allowbreak{}extendPoint} & 244-255 & helper &  \\
\texttt{unstack\_\allowbreak{}map}, \texttt{unstack\_\allowbreak{}eq\_\allowbreak{}of\_\allowbreak{}window\_\allowbreak{}eq} & 292, 304 & helper &  \\
\texttt{unstack\_\allowbreak{}eval} & 319 & helper & selection identity, one ring \\
\texttt{unstack\_\allowbreak{}eval₂\_\allowbreak{}eq\_\allowbreak{}sumCube} & 352 & helper (only \texttt{BlockClaims} uses it) & the identity as a cube sum \\
\end{longtable}}

\subsection{\file{LeanerVM/Protocol/ToCompPoly/AmbientStacking.lean}}

(generic; 10 public.)


\begin{leancode}
-- LeanerVM/Protocol/ToCompPoly/AmbientStacking.lean:69-70
def selectorWeight {μ : ℕ} (hμ : B.total ≤ 2 ^ μ) (b : Fin B.n) (z : Vector R μ) : R :=
  (lagrangeBasis (B.highPoint hμ b z))[B.selector hμ b]
\end{leancode}

\src{LeanerVM/Protocol/ToCompPoly/AmbientStacking.lean, at b435631; further ranges in the comments}

\texttt{selectorWeight b z = eq(\allowbreak{}sel\_\allowbreak{}b,\allowbreak{} z\_\allowbreak{}hi)}, the weight of block \texttt{b} at a stack point.

\par\medskip\noindent\textbf{Load-bearing and interface declarations.}

{\footnotesize
\begin{longtable}{@{}>{\raggedright\arraybackslash}p{0.280\linewidth}>{\raggedright\arraybackslash}p{0.080\linewidth}>{\raggedright\arraybackslash}p{0.260\linewidth}>{\raggedright\arraybackslash}p{0.300\linewidth}@{}}
\toprule
Declaration & Line & Class & Meaning \\
\midrule
\endfirsthead
\toprule
Declaration & Line & Class & Meaning \\
\midrule
\endhead
\bottomrule
\endfoot
\texttt{selectorWeight} & 69 & interface (in \texttt{stack\_\allowbreak{}eval\_\allowbreak{}ambient}; L once a verifier computes \texttt{eq(\allowbreak{}sel\_\allowbreak{}b,\allowbreak{} ζ\_\allowbreak{}hi)} with it) & above \\
\texttt{stack\_\allowbreak{}eval\_\allowbreak{}ambient} & 134 & interface (listed) & the stack at any point (below, B.3) \\
\end{longtable}}

\par\medskip\noindent\textbf{Helpers.}

{\footnotesize
\begin{longtable}{@{}>{\raggedright\arraybackslash}p{0.300\linewidth}>{\raggedright\arraybackslash}p{0.080\linewidth}>{\raggedright\arraybackslash}p{0.200\linewidth}>{\raggedright\arraybackslash}p{0.340\linewidth}@{}}
\toprule
Declaration & Line & Class & Meaning \\
\midrule
\endfirsthead
\toprule
Declaration & Line & Class & Meaning \\
\midrule
\endhead
\bottomrule
\endfoot
\texttt{selectorWeight\_\allowbreak{}extendPoint\_\allowbreak{}self} & 73 & helper & weight 1 at a lifted point of the block \\
\texttt{windowTable}, \texttt{windowTable\_\allowbreak{}getElem}, \texttt{evalMle\_\allowbreak{}windowTable} & 80, 86, 102 & helper & a table placed in one window \\
\texttt{stackAt\_\allowbreak{}decomposition} & 111 & helper & stack = pad + $Σ$ windows of (block $−$ pad) \\
\texttt{stack\_\allowbreak{}eval\_\allowbreak{}ambient\_\allowbreak{}zero} & 159 & helper & pad 0 \\
\texttt{sum\_\allowbreak{}selectorWeight\_\allowbreak{}of\_\allowbreak{}total\_\allowbreak{}eq} & 166 & helper & a full layout's weights sum to 1 \\
\texttt{stack\_\allowbreak{}eval₂\_\allowbreak{}ambient} & 178 & helper (used by \texttt{Stack.\allowbreak{}lean}) & at a point of another ring \\
\end{longtable}}

\subsection{\file{LeanerVM/Protocol/Stack.lean}}

(leanVM's: specialises to \texttt{K}, \texttt{E}, the spine's \texttt{Column} and \texttt{Layout}; 9 public.)


\begin{leancode}
-- LeanerVM/Protocol/Stack.lean:52-59
def Layout.comap {μ : ℕ} {ι ι' : Type} {κ : ι → ℕ} {κ' : ι' → ℕ} (L : Layout μ ι κ)
    (f : ι' → ι) (h : ∀ c, κ (f c) = κ' c) : Layout μ ι' κ' where
  read := fun q c ↦ ⟨Vector.cast (congrArg (2 ^ ·) (h c)) (L.read q (f c)).values⟩
  extend := fun c z ↦ L.extend (f c) (Vector.cast (h c).symm z)
  read_eval := fun q c z ↦ by … (3 lines)
-- :70, :73-74
def stackColumn (t : B.Tables K) (μ : ℕ) : Column μ := ⟨B.stackAt t μ 0⟩
def readColumn (hμ : B.total ≤ 2 ^ μ) (q : Column μ) (b : Fin B.n) : Column (B.size b) :=
  ⟨B.unstack hμ q.values b⟩
-- :97-100
def layout (hμ : B.total ≤ 2 ^ μ) : Layout μ (Fin B.n) B.size where
  read := B.readColumn hμ
  extend := B.extendPoint hμ
  read_eval := B.readColumn_eval hμ
\end{leancode}

\src{LeanerVM/Protocol/Stack.lean, at b435631; further ranges in the comments}


\par\medskip\noindent\textbf{Load-bearing and interface declarations.}

{\footnotesize
\begin{longtable}{@{}>{\raggedright\arraybackslash}p{0.280\linewidth}>{\raggedright\arraybackslash}p{0.080\linewidth}>{\raggedright\arraybackslash}p{0.260\linewidth}>{\raggedright\arraybackslash}p{0.300\linewidth}@{}}
\toprule
Declaration & Line & Class & Meaning \\
\midrule
\endfirsthead
\toprule
Declaration & Line & Class & Meaning \\
\midrule
\endhead
\bottomrule
\endfoot
\texttt{Layout.\allowbreak{}comap L f h} & 52 & load-bearing & column \texttt{c} of the new layout is column \texttt{f c} of \texttt{L} \\
\texttt{Blocks.\allowbreak{}stackColumn} & 70 & interface (listed; \texttt{stackOf}) & the witness stack, pad 0 \\
\texttt{Blocks.\allowbreak{}readColumn} & 73 & load-bearing & block \texttt{b} of a column \\
\texttt{Blocks.\allowbreak{}readColumn\_\allowbreak{}eval} & 78 & interface (listed) & the reading law, every \texttt{q} \\
\texttt{Blocks.\allowbreak{}readColumn\_\allowbreak{}stackColumn} & 85 & interface & reading the honest stack returns the block \\
\texttt{Blocks.\allowbreak{}layout} & 97 & load-bearing & the aligned blocks as a \texttt{Layout} \\
\texttt{Blocks.\allowbreak{}stackColumn\_\allowbreak{}eval\_\allowbreak{}ambient} & 104 & interface (not listed) & the zero-padded stack at any point \\
\texttt{Blocks.\allowbreak{}stack\_\allowbreak{}eval\_\allowbreak{}ambient\_\allowbreak{}one} & 118 & interface (listed; the leaf decomposition of Layer 6) & the one-padded stack at any point, char. 2 form \\
\end{longtable}}

\par\medskip\noindent\textbf{Helpers.}

{\footnotesize
\begin{longtable}{@{}>{\raggedright\arraybackslash}p{0.300\linewidth}>{\raggedright\arraybackslash}p{0.080\linewidth}>{\raggedright\arraybackslash}p{0.200\linewidth}>{\raggedright\arraybackslash}p{0.340\linewidth}@{}}
\toprule
Declaration & Line & Class & Meaning \\
\midrule
\endfirsthead
\toprule
Declaration & Line & Class & Meaning \\
\midrule
\endhead
\bottomrule
\endfoot
\texttt{Blocks.\allowbreak{}stackColumn\_\allowbreak{}eval} & 90 & helper & the honest stack at a lifted point \\
\end{longtable}}

\subsection{\file{LeanerVM/Protocol/Padding.lean}}

(leanVM's: the protocol's choice of slice, over any ring; 6 public.)


\begin{leancode}
-- LeanerVM/Protocol/Padding.lean:43-44, 60-61
def prodVars (m : ℕ) : CMlPolynomialEval R m :=
  Vector.ofFn fun i ↦ if i = onesIndex m then 1 else 0
def padHigh {k : ℕ} (t : CMlPolynomialEval R k) (m : ℕ) : CMlPolynomialEval R (k + m) :=
  placeSlice t (onesIndex m)
\end{leancode}

\src{LeanerVM/Protocol/Padding.lean, at b435631; further ranges in the comments}


\par\medskip\noindent\textbf{Load-bearing and interface declarations.}

{\footnotesize
\begin{longtable}{@{}>{\raggedright\arraybackslash}p{0.280\linewidth}>{\raggedright\arraybackslash}p{0.080\linewidth}>{\raggedright\arraybackslash}p{0.260\linewidth}>{\raggedright\arraybackslash}p{0.300\linewidth}@{}}
\toprule
Declaration & Line & Class & Meaning \\
\midrule
\endfirsthead
\toprule
Declaration & Line & Class & Meaning \\
\midrule
\endhead
\bottomrule
\endfoot
\texttt{prodVars m} & 43 & interface (listed) & the table of \texttt{x\_\allowbreak{}0 ⋯ x\_\allowbreak{}\{m−1\}} \\
\texttt{sumCube\_\allowbreak{}prodVars} & 47 & interface (listed; acceptance test 7) & \texttt{Σ\_\allowbreak{}x x\_\allowbreak{}0⋯x\_\allowbreak{}\{m−1\} = 1} \\
\texttt{padHigh t m} & 60 & interface (listed) & \texttt{t} lifted by the product of \texttt{m} new high variables \\
\texttt{sumCube\_\allowbreak{}padHigh} & 64 & interface (listed) & keeps the cube sum (completeness of the table sumcheck) \\
\texttt{evalMle\_\allowbreak{}padHigh} & 70 & interface (\textbf{not listed}; the verifier's final check needs it) & extension \texttt{= t̃(\allowbreak{}z)·∏\_\allowbreak{}b s\_\allowbreak{}b} \\
\end{longtable}}

\par\medskip\noindent\textbf{Helpers.}

{\footnotesize
\begin{longtable}{@{}>{\raggedright\arraybackslash}p{0.300\linewidth}>{\raggedright\arraybackslash}p{0.080\linewidth}>{\raggedright\arraybackslash}p{0.200\linewidth}>{\raggedright\arraybackslash}p{0.340\linewidth}@{}}
\toprule
Declaration & Line & Class & Meaning \\
\midrule
\endfirsthead
\toprule
Declaration & Line & Class & Meaning \\
\midrule
\endhead
\bottomrule
\endfoot
\texttt{evalMle\_\allowbreak{}prodVars} & 51 & helper & its extension is \texttt{∏ s\_\allowbreak{}b} \\
\end{longtable}}

\subsection{\file{LeanerVM/Protocol/ClaimWeights.lean}}

(leanVM's: over an abstract instance; 4 public.)


\begin{leancode}
-- LeanerVM/Protocol/ClaimWeights.lean:45-48
def eqWeight {μ : ℕ} (p : Vector E μ) : Weight μ where
  onCube := lagrangeBasis p
  mle := fun r ↦ ∏ k : Fin μ, ((1 - r[k]) * (1 - p[k]) + r[k] * p[k])
  mle_eq := fun r ↦ (evalMle_lagrangeBasis p r).symm
\end{leancode}

\src{LeanerVM/Protocol/ClaimWeights.lean, at b435631; further ranges in the comments}


\par\medskip\noindent\textbf{Load-bearing and interface declarations.}

{\footnotesize
\begin{longtable}{@{}>{\raggedright\arraybackslash}p{0.280\linewidth}>{\raggedright\arraybackslash}p{0.080\linewidth}>{\raggedright\arraybackslash}p{0.260\linewidth}>{\raggedright\arraybackslash}p{0.300\linewidth}@{}}
\toprule
Declaration & Line & Class & Meaning \\
\midrule
\endfirsthead
\toprule
Declaration & Line & Class & Meaning \\
\midrule
\endhead
\bottomrule
\endfoot
\texttt{eqWeight p} & 45 & load-bearing (the opening phase's weights; the verifier runs \texttt{mle}) & the weight \texttt{eq(\allowbreak{}p,\allowbreak{} ·)} \\
\texttt{ColumnClaim.\allowbreak{}holds\_\allowbreak{}iff\_\allowbreak{}weighted} & 58 & interface (listed; the opening phase) & a column claim is the weighted claim with weight \texttt{eq(\allowbreak{}extend c z,\allowbreak{} ·)} \\
\end{longtable}}

\par\medskip\noindent\textbf{Helpers.}

{\footnotesize
\begin{longtable}{@{}>{\raggedright\arraybackslash}p{0.300\linewidth}>{\raggedright\arraybackslash}p{0.080\linewidth}>{\raggedright\arraybackslash}p{0.200\linewidth}>{\raggedright\arraybackslash}p{0.340\linewidth}@{}}
\toprule
Declaration & Line & Class & Meaning \\
\midrule
\endfirsthead
\toprule
Declaration & Line & Class & Meaning \\
\midrule
\endhead
\bottomrule
\endfoot
\texttt{Weight.\allowbreak{}pair\_\allowbreak{}eq\_\allowbreak{}sumCube} & 34 & helper & the pairing as \texttt{sumCube (\allowbreak{}hadamard …)} \\
\texttt{eqWeight\_\allowbreak{}pair} & 51 & helper & pairing \texttt{eq(\allowbreak{}p,\allowbreak{} ·)} with \texttt{q} is \texttt{q̃(\allowbreak{}p)} \\
\end{longtable}}

\subsection{\file{LeanerVM/Protocol/BlockClaims.lean}}

(generic content, placed in the leanVM half; 8 public; no consumer.)


\par\medskip\noindent\textbf{Helpers.}

{\footnotesize
\begin{longtable}{@{}>{\raggedright\arraybackslash}p{0.300\linewidth}>{\raggedright\arraybackslash}p{0.080\linewidth}>{\raggedright\arraybackslash}p{0.200\linewidth}>{\raggedright\arraybackslash}p{0.340\linewidth}@{}}
\toprule
Declaration & Line & Class & Meaning \\
\midrule
\endfirsthead
\toprule
Declaration & Line & Class & Meaning \\
\midrule
\endhead
\bottomrule
\endfoot
\texttt{BlockClaim B S} & 54 & helper (unused) & (block, point, value) \\
\texttt{ambientPoint}, \texttt{weight} & 69, 73 & helper (unused) & the lifted point; \texttt{lagrangeBasis} of it \\
\texttt{IsValid φ hμ q} & 78 & helper (unused) & \texttt{eval₂Mle (\allowbreak{}unstack q block) φ point = value} \\
\texttt{pairing\_\allowbreak{}eq}, \texttt{isValid\_\allowbreak{}iff\_\allowbreak{}pairing} & 83, 90 & helper (unused) & \texttt{unstack\_\allowbreak{}eval₂\_\allowbreak{}eq\_\allowbreak{}sumCube} restated \\
\texttt{isValid\_\allowbreak{}iff\_\allowbreak{}of\_\allowbreak{}map\_\allowbreak{}window\_\allowbreak{}eq}, \texttt{isValid\_\allowbreak{}iff\_\allowbreak{}of\_\allowbreak{}window\_\allowbreak{}eq} & 101, 113 & helper (unused) & locality: only the window matters \\
\end{longtable}}

\texttt{git grep -{}w BlockClaim b435631 -{}-{} LeanerVM tests} finds \texttt{BlockClaims.\allowbreak{}lean} and its test only.

\subsection{\file{LeanerVM/Protocol/FixedColumns.lean}}

(leanVM's; imports the arithmetization; 12 public.)


\begin{leancode}
-- LeanerVM/Protocol/FixedColumns.lean:48, 56-57
def idxColumn (κ : ℕ) : Column κ := ⟨powersTable g κ⟩
def idxColumnEval {κ : ℕ} (z : Vector E κ) : E :=
  ∏ j : Fin κ, ((1 - z[j]) + z[j] * algebraMap K E (g ^ (2 ^ j.val)))
-- :79-82
def bytecodeColumn (prog : Program) : Column (prog.logSize + 4) :=
  ⟨Vector.ofFn fun i ↦
    let p := (cubeSplit prog.logSize 4).symm i
    (encodeSlots (prog.code p.1))[p.2]⟩
-- :108-111
def bytecodeColumnEval (prog : Program) (z : Vector E prog.logSize) (w : Vector E 4) : E :=
  ∑ s : Fin 16, (lagrangeBasis w)[s] *
    ∑ i : Fin (2 ^ prog.logSize),
      algebraMap K E ((encodeSlots (prog.code i))[s]) * (lagrangeBasis z)[i]
-- LeanerVM/Arithmetization/Bytecode.lean:94-96 (what bytecodeColumn puts in a cell)
def encodeSlots (i : Instr) : Vector K 16 :=
  let e := entry i
  #v[0, 0, 0, e[0], e[1], e[2], e[3], e[4], e[5], e[6], e[7], 0, 0, 0, 0, 0]
\end{leancode}

\src{LeanerVM/Protocol/FixedColumns.lean, at b435631; further ranges in the comments}


\par\medskip\noindent\textbf{Load-bearing and interface declarations.}

{\footnotesize
\begin{longtable}{@{}>{\raggedright\arraybackslash}p{0.280\linewidth}>{\raggedright\arraybackslash}p{0.080\linewidth}>{\raggedright\arraybackslash}p{0.260\linewidth}>{\raggedright\arraybackslash}p{0.300\linewidth}@{}}
\toprule
Declaration & Line & Class & Meaning \\
\midrule
\endfirsthead
\toprule
Declaration & Line & Class & Meaning \\
\midrule
\endhead
\bottomrule
\endfoot
\texttt{idxColumn κ} & 48 & load-bearing (a \texttt{Coord.\allowbreak{}known} column of the instance) & cell \texttt{i} holds \texttt{g\textasciicircum{}i} \\
\texttt{idxColumnEval} & 56 & load-bearing (the verifier) & \texttt{∏\_\allowbreak{}j (\allowbreak{}(\allowbreak{}1 − z\_\allowbreak{}j) + z\_\allowbreak{}j g\textasciicircum{}\{2\textasciicircum{}j\})} \\
\texttt{idxColumn\_\allowbreak{}eval} & 60 & interface (listed) & the oracle's answer is \texttt{idxColumnEval} \\
\texttt{idxColumnEval\_\allowbreak{}eq} & 66 & interface (listed) & the specification's form \texttt{∏ (\allowbreak{}1 + z\_\allowbreak{}k(\allowbreak{}1 + g\textasciicircum{}\{2\textasciicircum{}k\}))} \\
\texttt{bytecodeColumn prog} & 79 & load-bearing (a \texttt{Coord.\allowbreak{}known} column; the verifier) & cell \texttt{i + 2\textasciicircum{}κ·s} holds slot \texttt{s} of instruction \texttt{i} \\
\texttt{bytecodeColumn\_\allowbreak{}answer\_\allowbreak{}boolVec} & 97 & interface (listed) & at the cube point \texttt{(\allowbreak{}bits of i,\allowbreak{} bits of s)} the oracle answers slot \texttt{s} of instruction \texttt{i} \\
\texttt{bytecodeColumnEval} & 108 & load-bearing (the verifier) & \texttt{Σ\_\allowbreak{}s eq(\allowbreak{}w,\allowbreak{} s) Σ\_\allowbreak{}i slot\_\allowbreak{}s(\allowbreak{}i) eq(\allowbreak{}z,\allowbreak{} i)} \\
\texttt{bytecodeColumn\_\allowbreak{}eval} & 114 & interface (listed) & the oracle's answer at \texttt{z ++ w} is \texttt{bytecodeColumnEval prog z w} \\
\end{longtable}}

\par\medskip\noindent\textbf{Helpers.}

{\footnotesize
\begin{longtable}{@{}>{\raggedright\arraybackslash}p{0.300\linewidth}>{\raggedright\arraybackslash}p{0.080\linewidth}>{\raggedright\arraybackslash}p{0.200\linewidth}>{\raggedright\arraybackslash}p{0.340\linewidth}@{}}
\toprule
Declaration & Line & Class & Meaning \\
\midrule
\endfirsthead
\toprule
Declaration & Line & Class & Meaning \\
\midrule
\endhead
\bottomrule
\endfoot
\texttt{idxColumn\_\allowbreak{}get} & 51 & helper & the cells \\
\texttt{bytecodeSlotColumn prog s} & 75 & helper & slot \texttt{s} of every instruction \\
\texttt{bytecodeColumn\_\allowbreak{}slot} & 85 & helper (restates the definition; section C) &  \\
\texttt{slice\_\allowbreak{}bytecodeColumn} & 90 & helper & slice \texttt{s} is \texttt{bytecodeSlotColumn s} \\
\end{longtable}}

\textbf{Tally.} 143 public declarations: load-bearing 20 (\texttt{cubeIndex}, \texttt{cubeSplit}, \texttt{boolVec}, \texttt{slice}, \texttt{powersTable}, \texttt{Blocks}, \texttt{offsetNat}, \texttt{offset}, \texttt{total}, \texttt{selector}, \texttt{extendPoint}, \texttt{unstack}, \texttt{Layout.\allowbreak{}comap}, \texttt{readColumn}, \texttt{layout}, \texttt{eqWeight}, \texttt{idxColumn}, \texttt{idxColumnEval}, \texttt{bytecodeColumn}, \texttt{bytecodeColumnEval}); interface 41; helper 82 (of which 8 in \texttt{BlockClaims.\allowbreak{}lean}, unused). Per module (L / I / H): Multilinear 4 / 10 / 29, BitProductTable 1 / 4 / 8, Stacking 7 / 10 / 21, AmbientStacking 0 / 2 / 8, Stack 3 / 5 / 1, Padding 0 / 5 / 1, ClaimWeights 1 / 1 / 2, BlockClaims 0 / 0 / 8, FixedColumns 4 / 4 / 4. Generic modules: the four under \texttt{ToCompPoly/} (104 declarations); \texttt{BlockClaims.\allowbreak{}lean} is generic in content too (any \texttt{Blocks}, any rings) but sits in the leanVM half.
